% !TEX root = ../Main.tex
\chapter{概率（二）：随机事件与古典概型}

\section{随机事件}

\subsection*{事件的运算}

\begin{itemize}
    \item \textbf{交事件（积事件）}：$A$ 与 $B$ \textbf{同时}发生，记作 $A\cap B$（或 $AB$）；
    \item \textbf{并事件（和事件）}：$A$、$B$ \textbf{至少一个}发生，记作 $A\cup B$。
\end{itemize}

\subsection*{事件的关系}

\begin{itemize}
    \item \textbf{包含}：$A$ 发生则 $B$ 一定发生，记 $A\subseteq B$。特别地，$A\subseteq B$ 且 $B\subseteq A$ 时 $A$、$B$ \textbf{相等}；
    \item \textbf{互斥}：$A$、$B$ 不能同时发生，$A\cap B=\varnothing$。特别地，若还"必有一个发生"，就是\textbf{对立}（非此即彼），$B=\bar{A}$；
    \item \textbf{独立}：$A$ 发生与否不影响 $B$ 的概率，判定式是 $P(AB)=P(A)P(B)$。
\end{itemize}

\section{古典概型}

\state{古典概型的两个前提}{
    \begin{itemize}
        \item \textbf{有限性}：基本事件只有有限个；
        \item \textbf{等可能性}：每个基本事件发生的可能性相等。
    \end{itemize}
}

满足这两条时，事件 $A$ 的概率为
\[
P(A)=\frac{A\ \text{包含的基本事件数}}{\text{总基本事件数}}.
\]

所以做古典概型，核心就是\textbf{把基本事件数清楚}——常用的就是\textbf{穷举法}：把所有情形一个不漏地列出来数。

\section{两道古典概型例题}

\ex[例 1]{在 $0,1,2,3,4,5,6,7,8$ 这 $9$ 个数中\textbf{不放回}地取出两个。记事件 $A$ 为"取出的两个数中有且仅有一个为奇数"，事件 $B$ 为"取出的两个数中有且仅有一个为偶数"。问是否有 $P(A)=P(B)$？}

\sol{
    这 $9$ 个数里奇数有 $1,3,5,7$ 共 $4$ 个，偶数有 $0,2,4,6,8$ 共 $5$ 个。不放回取两个、不计顺序，总基本事件数 $C_9^2=36$。

    事件 $A$"恰好一个奇数"就是\textbf{一奇一偶}：
    \[
    P(A)=\frac{C_4^1\,C_5^1}{C_9^2}=\frac{4\times 5}{36}=\frac{20}{36}=\frac{5}{9}.
    \]
    取两个数"恰好一个奇数"和"恰好一个偶数"说的是同一件事——一奇一偶，所以 $A$、$B$ \textbf{本来就是同一个事件}，自然
    \[
    P(A)=P(B)=\frac{5}{9}.
    \]
}

\ex[例 1 变形]{把"不放回"改成\textbf{有放回}，其它条件不变，问是否还有 $P(A)=P(B)$？}

\sol{
    有放回取两次、分先后，总基本事件数 $9\times 9=81$。

    事件 $A$"恰好一次取到奇数"：哪一次取奇数有 $2$ 种，奇数 $4$ 个、偶数 $5$ 个，
    \[
    P(A)=\frac{2\times 4\times 5}{81}=\frac{40}{81}.
    \]
    同理 $B$"恰好一次取到偶数"也是恰好一次取奇数一次取偶数，$P(B)=\dfrac{2\times 5\times 4}{81}=\dfrac{40}{81}$。仍然
    \[
    P(A)=P(B)=\frac{40}{81}.
    \]
    取两次时"恰好一个奇数"与"恰好一个偶数"永远是同一件事，所以放不放回都相等。
}

\ex[例 2]{一只器皿中装有 $7$ 个红玻璃球、$3$ 个绿玻璃球。从中\textbf{无放回}地任取 $2$ 次，每次取 $1$ 个。求：取得两个红球的概率、取得两个绿球的概率、取得两个同色球的概率、至少取得一个红球的概率。}

\sol{
    共 $10$ 个球，无放回取两个、不计顺序，总基本事件数 $C_{10}^2=45$。

    取得两红：
    \[
    P(\text{两红})=\frac{C_7^2}{C_{10}^2}=\frac{21}{45}=\frac{7}{15}.
    \]
    取得两绿：
    \[
    P(\text{两绿})=\frac{C_3^2}{C_{10}^2}=\frac{3}{45}=\frac{1}{15}.
    \]
    "两个同色"就是"两红"或"两绿"，这两个事件互斥，概率相加：
    \[
    P(\text{同色})=\frac{7}{15}+\frac{1}{15}=\frac{8}{15}.
    \]
    "至少一个红"的对立事件是"一个红都没有"，即"两个都是绿"，用对立事件最省事：
    \[
    P(\text{至少一红})=1-P(\text{两绿})=1-\frac{1}{15}=\frac{14}{15}.
    \]
}
